\thispagestyle{plain}
\begin{center}
{\small\sffamily\color{accent}ECONOMIC MEASUREMENT AND TEXT AS DATA}\par\vspace{11pt}
{\LARGE\bfseries\color{ink}Low-Cost Extraction of Economic\\[3pt] Forecasts from Newspapers}\par\vspace{9pt}
{\large Human Validation, Measurement Sensitivity,\\and Predictive Applications}\par\vspace{12pt}
{\small Research manuscript\enspace$\vert$\enspace Revised October 22, 2026}
\end{center}
\vspace{3pt}
\begin{abstract}
Newspapers preserve economic expectations whose target, timing, and attribution must be interpreted before a forecast can be evaluated. This paper studies inexpensive structured extraction using ProQuest exports for 1900--2010 and a separate exploratory extension for 2011--2026. The historical corpus contains 274,070 machine-verified candidates, with 175,855 scored under deterministic rules; the modern extension contains 91,118 retained candidates, with 56,886 mature, supported claims evaluated. Separate Library of Congress development experiments project API costs of \$26.82 versus \$403.70 for 23,000 pages, with respective F1 scores of 0.636 and 0.784 against a 52-claim model-adjudicated reference. These are development comparisons, not ProQuest production costs or human quality estimates. Two human reviewers evaluated the same 300 sampled candidates: 250 historical and 50 modern. Their decisions are preserved separately; inter-reviewer agreement has not yet been calculated. Existing first-review field comparisons reveal weaknesses in horizons and geographic scope. The central historical result is measurement sensitivity: on 9,260 matched employment claims, article-weighted accuracy ranges from 34.2\% to 57.6\% across outcome bands, while changing the employment benchmark reverses the baseline-relative point estimate. In the modern extension, a frozen recurrent model achieves AUC 0.738 versus 0.654 for a metadata-plus-history logistic control, but paired block intervals do not establish the Brier-score improvement and employment AUC is only 0.306. Modern results remain exploratory because prior outcome exposure is undocumented. The evidence supports inexpensive extraction as candidate generation whose economic use requires human source review, explicit targets, matched baselines, and chronological evaluation.
\end{abstract}
\noindent{\small\textbf{Keywords:} historical expectations; newspaper forecasts; language models; text measurement; forecast evaluation.}
\vspace{4pt}

